% GENERATED FILE. Edit canonical lessons, then run npm run generate:books.
% canonical-source-hash: fnv1a64:f4d520c7d70036f0
% canonical-lessons: AR-C84-farigh, AR-C84-ghani, AR-C84-faqir, AR-C84-nazif, AR-C84-wasikh

\chapter{Empty, Rich, Poor, Clean, Dirty}
\label{ch:ar-empty-rich-poor-clean-dirty}
\hlchaptermodality{84}

\begin{chapteropening}
\textbf{By the end of this chapter:} \emph{I can say empty, rich, poor, clean and dirty.}

Complete the last lesson of chapter 84: I can say empty, rich, poor, clean and dirty.
\end{chapteropening}

\section[\texorpdfstring{\ar{فارغ} (fārigh)}{فارغ (fārigh)}]{\ar{فارغ} (fārigh) — empty}

\label{lesson:AR-C84-farigh}



\begin{quote}
Before the new one: say the Arabic for difficult, then the Arabic for cheap.
\end{quote}

\subsection*{You'll want to know: \ar{فارغ}}
\textbf{\ar{فارغ}} — \emph{fārigh} — \textquotedblleft{}empty\textquotedblright{}.

\begin{cousinweb}
\textbf{\ar{فارغ}} (\emph{fārigh}) is built on the native Arabic root \emph{f-r-gh}, \textquotedblleft{}to be empty, to be free\textquotedblright{} (Lane). \emph{Fārigh} is the active participle, \textquotedblleft{}empty, vacant\textquotedblright{}.
\end{cousinweb}

\begin{grammarlens}[title={What you've built}]
One more word for this chapter.
\end{grammarlens}

\subsection*{Your turn}
\begin{itemize}
\raggedright
  \item \emph{Say it:} \emph{fārigh}
  \item \emph{Say it:} \emph{fārigh}, once more
  \item \emph{Say it:} say \emph{fārigh}
  \item \emph{Recall:} say the Arabic for difficult, then the Arabic for cheap, then say \emph{fārigh} again
\end{itemize}

\begin{tcolorbox}[breakable,colback=teal!4,colframe=teal!35!black,title={Before you move on}]
What does \ar{فارغ} mean? (\textquotedblleft{}Empty\textquotedblright{}.) Say it once more.
\end{tcolorbox}
\section[\texorpdfstring{\ar{غني} (ghanī)}{غني (ghanī)}]{\ar{غني} (ghanī) — rich}

\label{lesson:AR-C84-ghani}



\begin{quote}
Before the new one: say the Arabic for cheap, then the Arabic for empty.
\end{quote}

\subsection*{You'll want to know: \ar{غني}}
\textbf{\ar{غني}} — \emph{ghanī} — \textquotedblleft{}rich\textquotedblright{}.

\begin{cousinweb}
\textbf{\ar{غني}} (\emph{ghanī}) is built on the native Arabic root \emph{gh-n-y}, \textquotedblleft{}to be free from want, to be rich\textquotedblright{} (Lane). \emph{Ghanī} is the \emph{faʿīl} adjective, \textquotedblleft{}rich\textquotedblright{}.
\end{cousinweb}

\begin{grammarlens}[title={What you've built}]
One more word for this chapter.
\end{grammarlens}

\subsection*{Your turn}
\begin{itemize}
\raggedright
  \item \emph{Say it:} \emph{ghanī}
  \item \emph{Say it:} \emph{ghanī}, once more
  \item \emph{Say it:} say \emph{ghanī}
  \item \emph{Recall:} say the Arabic for cheap, then the Arabic for empty, then say \emph{ghanī} again
\end{itemize}

\begin{tcolorbox}[breakable,colback=teal!4,colframe=teal!35!black,title={Before you move on}]
What does \ar{غني} mean? (\textquotedblleft{}Rich\textquotedblright{}.) Say it once more.
\end{tcolorbox}
\section[\texorpdfstring{\ar{فقير} (faqīr)}{فقير (faqīr)}]{\ar{فقير} (faqīr) — poor}

\label{lesson:AR-C84-faqir}



\begin{quote}
Before the new one: say the Arabic for empty, then the Arabic for rich.
\end{quote}

\subsection*{You'll want to know: \ar{فقير}}
\textbf{\ar{فقير}} — \emph{faqīr} — \textquotedblleft{}poor\textquotedblright{}.

\begin{cousinweb}
\textbf{\ar{فقير}} (\emph{faqīr}) is built on the native Arabic root \emph{f-q-r}, \textquotedblleft{}to be poor; a vertebra\textquotedblright{} (Lane). English \emph{fakir} comes from \emph{faqīr}.
\end{cousinweb}

\begin{grammarlens}[title={What you've built}]
One more word for this chapter.
\end{grammarlens}

\subsection*{Your turn}
\begin{itemize}
\raggedright
  \item \emph{Say it:} \emph{faqīr}
  \item \emph{Say it:} \emph{faqīr}, once more
  \item \emph{Say it:} say \emph{faqīr}
  \item \emph{Recall:} say the Arabic for empty, then the Arabic for rich, then say \emph{faqīr} again
\end{itemize}

\begin{tcolorbox}[breakable,colback=teal!4,colframe=teal!35!black,title={Before you move on}]
What does \ar{فقير} mean? (\textquotedblleft{}Poor\textquotedblright{}.) Say it once more.
\end{tcolorbox}
\section[\texorpdfstring{\ar{نظيف} (naẓīf)}{نظيف (naẓīf)}]{\ar{نظيف} (naẓīf) — clean}

\label{lesson:AR-C84-nazif}



\begin{quote}
Before the new one: say the Arabic for rich, then the Arabic for poor.
\end{quote}

\subsection*{You'll want to know: \ar{نظيف}}
\textbf{\ar{نظيف}} — \emph{naẓīf} — \textquotedblleft{}clean\textquotedblright{}.

\begin{cousinweb}
\textbf{\ar{نظيف}} (\emph{naẓīf}) is built on the native Arabic root \emph{n-ẓ-f}, \textquotedblleft{}to be clean\textquotedblright{}; Lane's short article on this root does not treat it. \emph{Naẓīf} is the \emph{faʿīl} adjective, \textquotedblleft{}clean\textquotedblright{}.
\end{cousinweb}

\begin{grammarlens}[title={What you've built}]
One more word for this chapter.
\end{grammarlens}

\subsection*{Your turn}
\begin{itemize}
\raggedright
  \item \emph{Say it:} \emph{naẓīf}
  \item \emph{Say it:} \emph{naẓīf}, once more
  \item \emph{Say it:} say \emph{naẓīf}
  \item \emph{Recall:} say the Arabic for rich, then the Arabic for poor, then say \emph{naẓīf} again
\end{itemize}

\begin{tcolorbox}[breakable,colback=teal!4,colframe=teal!35!black,title={Before you move on}]
What does \ar{نظيف} mean? (\textquotedblleft{}Clean\textquotedblright{}.) Say it once more.
\end{tcolorbox}
\section[\texorpdfstring{\ar{وسخ} (wasikh)}{وسخ (wasikh)}]{\ar{وسخ} (wasikh) — dirty}

\label{lesson:AR-C84-wasikh}



\begin{quote}
Before the new one: say the Arabic for poor, then the Arabic for clean.
\end{quote}

\subsection*{You'll want to know: \ar{وسخ}}
\textbf{\ar{وسخ}} — \emph{wasikh} — \textquotedblleft{}dirty\textquotedblright{}.

\begin{cousinweb}
\textbf{\ar{وسخ}} (\emph{wasikh}) is built on the native Arabic root \emph{w-s-kh}, \textquotedblleft{}to be dirty\textquotedblright{} (Lane). \emph{Wasikh} is the adjective, \textquotedblleft{}dirty\textquotedblright{}.
\end{cousinweb}

\begin{grammarlens}[title={What you've built}]
One more word for this chapter.
\end{grammarlens}

\subsection*{Your turn}
\begin{itemize}
\raggedright
  \item \emph{Say it:} \emph{wasikh}
  \item \emph{Say it:} \emph{wasikh}, once more
  \item \emph{Say it:} say \emph{wasikh}
  \item \emph{Recall:} say the Arabic for poor, then the Arabic for clean, then say \emph{wasikh} again
\end{itemize}

\begin{tcolorbox}[breakable,colback=teal!4,colframe=teal!35!black,title={Before you move on}]
What does \ar{وسخ} mean? (\textquotedblleft{}Dirty\textquotedblright{}.) Say it once more.
\end{tcolorbox}
